\section{Introduction} 
\label{sec:intro}
\subsection{Background}
\label{subsec:background}
Approximately 600 million people in sub-Saharan Africa (SSA) lack access to electricity,~\cite{iea2023} representing more than half of the global population without power. This energy access deficit creates significant barriers to socio-economic development, exacerbates environmental degradation, and contributes to climate change.~\cite{urban2019,esmap2019}

Energy poverty in rural SSA manifests through multiple interconnected dimensions. Health complications arise from indoor air pollution caused by reliance on traditional biomass for cooking and kerosene for lighting.~\cite{lam2012} Economic opportunities remain constrained by inadequate infrastructure to support income-generating activities and small businesses.~\cite{bos2018} Educational outcomes suffer as children lack reliable lighting for evening study, while limited access to communication technologies perpetuates cycles of poverty and underdevelopment.~\cite{kangawa2008}

Solar mini-grids have emerged as technically viable decentralized solutions for communities too remote for cost-effective grid connection yet sufficiently populated to support standalone commercial operations.~\cite{se4all2015} Solar mini-grids effect significant socio-economic benefits in recipient communities,~\cite{minigrid-impact} and the United Nations estimates decentralized renewable energy systems will provide clean energy access to two-thirds of rural SSA populations.~\cite{irena2018} Despite this potential, mini-grid deployment remains constrained by high upfront capital costs, limited long-term financing, revenue uncertainty from rural consumers' constrained ability to pay, regulatory ambiguities, and high transaction costs relative to project size.~\cite{bhattacharyya2016,trotter2017,esmap2019}

Carbon markets have been proposed as supplementary financing mechanisms to address these constraints. By generating revenue from verified emission reductions, carbon credits could theoretically improve mini-grid commercial viability while simultaneously advancing climate mitigation objectives.~\cite{creutzig2017} However, post-Paris governance fragmentation, with competing standards, divergent accounting rules, and uncertain host government authorization under Article 6, has complicated implementation.~\cite{ahonen-governance} Recent events underscore these risks: KOKO Networks, serving 1.5 million Kenyan households, entered administration in February 2026 after the Kenyan government refused a Letter of Authorisation, illustrating the extreme fragility of business models dependent on carbon revenues.~\cite{koko-disrupt} Whether carbon markets can meaningfully contribute to scaling rural electrification in SSA requires empirical examination.

\subsection{Problem Statement}
\label{subsec:problem}
Significant knowledge gaps exist regarding the practical implementation and effectiveness of carbon finance for mini-grid development. Current literature provides limited empirical evidence on how carbon finance mechanisms perform in rural African contexts, what conditions enable successful integration, and how transaction costs compare to potential revenues at sub-200-kW scales. The complex regulatory landscape, including host government authorization requirements,~\cite{schneider-integrity,ahonen-governance} revenue leakage to intermediaries,~\cite{power-africa-credits} and undefined treatment of carbon revenues within tariff frameworks,~\cite{kreibich2021} creates uncertainty for developers and investors. Without clear understanding of these conditions, the potential of carbon markets to scale rural electrification remains largely theoretical rather than practical.

\subsection{Research Questions}
\label{subsec:rqs}
This research examines the potential of carbon markets in scaling up rural solar mini-grids in SSA addressing two questions:
\begin{enumerate}
    \item[RQ1:] How effectively do carbon finance mechanisms incentivize investment and improve commercial viability of rural solar mini-grid projects?
    \item[RQ2:] What technical, financial, regulatory, and policy conditions are necessary for successful carbon market integration at scale?
\end{enumerate}

\subsection{Significance and Scope}
\label{subsec:scope}
This research contributes empirical evidence on carbon market integration with rural electrification, advancing understanding of how market-based mechanisms can address development financing challenges. The findings inform policy frameworks for carbon market regulation, rural electrification strategies, actionable insights for developers, investors, and policy makers navigating post-Paris governance complexity. For developers, this research provides insights on project design and financing structures that enhance commercial viability through carbon market participation.

The research focuses on rural solar mini-grids in SSA, examining integration with voluntary carbon markets across East and West African contexts (i.e. Kenya and Nigeria) to capture variation in regional maturity and deployment models. Projects spanning small (<100 kWp), medium (100–400 kWp), and large (>1000 kWp) scales are examined to assess how carbon market viability varies by project size. The temporal scope focuses on developments from 2015 onwards, with particular attention to post-Paris Agreement development.

The findings are primarily applicable to SSA countries with foundational regulatory frameworks, adequate solar resources, and relatively stable political environments. Applicability beyond SSA depends on the presence of similar socio-economic, institutional, and development conditions.